\documentclass[11pt]{article}
\usepackage[ngerman]{babel}
\usepackage[a4paper,margin=0.9in]{geometry}
\usepackage[T1]{fontenc}
\usepackage[utf8]{inputenc}
\usepackage{lmodern}
\usepackage{amsmath,amssymb,amsthm,mathtools,mathrsfs}
\usepackage{microtype}
\usepackage{fancyhdr}
\usepackage{array,longtable,enumitem,multicol}
\usepackage{tikz}
\providecommand{\dd}{\,d}
\providecommand{\tht}{\vartheta}
\providecommand{\sn}{\operatorname{sn}}
\providecommand{\cn}{\operatorname{cn}}
\providecommand{\dn}{\operatorname{dn}}
\providecommand{\ctg}{\operatorname{cotg}}
\providecommand{\ii}{\mathrm{i}}
\providecommand{\Kp}{K'}
\providecommand{\wpfunc}{\wp}
\setlength{\parindent}{1.5em}
\setlength{\parskip}{0.18em}
\emergencystretch=6em
\pagestyle{fancy}
\fancyhf{}
\cfoot{\thepage}
\lhead{Heinrich Weber, Lehrbuch der Algebra, III. Band}
\rhead{Deutsche Transkription}
\begin{document}

\begin{center}
{\large Fünfter Abschnitt.}\\[0.7em]
{\Large Die Modulfunktionen.}
\end{center}

\section*{\S\,51. Die elliptischen Differentialgleichungen.}

Die bisher definierten elliptischen Funktionen sind Funktionen der beiden Variablen $v,\omega$, und die Moduln $\varkappa,\varkappa'$, ferner $j(\omega),K,K'$ sind von dem Periodenverhältnis $\omega$, das einen positiv imaginären Bestandteil haben muß, abhängig. Ebenso ist die Funktion $\wp(u)$ eine Funktion der drei Variablen $u,\omega_1,\omega_2$ und die Invarianten $g_2,g_3$ sind von $\omega_1,\omega_2$ abhängig. Die Aufgabe der Umkehrung der elliptischen Integrale, wie wir sie in \S\,14 formuliert haben, setzt aber voraus, daß in den elliptischen Funktionen $\varkappa^2$ (oder bei der Weierstrassschen Normalform $g_2,g_3$) beliebig gegebene Werte haben, und die vollständige Lösung dieser Aufgabe verlangt also, daß nicht $\omega$, sondern $\varkappa^2$ (bzw. $g_2,g_3$) als zweite unabhängige Variable auftritt, und daß $\omega$ durch diese ausgedrückt werde. Dieser Aufgabe werden die nächsten Betrachtungen gewidmet sein.

Wir gehen aus von der Aufgabe, ein System von Differentialgleichungen, welches wir das der elliptischen Differentialgleichungen nennen, zu integrieren:
\begin{equation}
\frac{\dd x}{\dd v}=yz,\qquad
\frac{\dd y}{\dd v}=-zx,\qquad
\frac{\dd z}{\dd v}=-\varkappa^2xy,
\tag{1}
\end{equation}
unter den Nebenbedingungen:
\begin{equation}
\text{für }v=0\text{ soll }x=0,
\quad y=1,
\quad z=1\text{ sein.}
\tag{2}
\end{equation}
Wir haben im vorigen Abschnitt gesehen, daß, wenn
\begin{equation}
\varkappa^2={\tht_{10}^4\over\tht_{00}^4}
\tag{3}
\end{equation}
ist, diese Aufgabe durch die elliptischen Funktionen gelöst wird, und zwar in der Weise, daß $x,y,z$ in der ganzen $v$-Ebene eindeutige und, außer wo sie unendlich sind, stetige Funktionen von $v$ werden. Ist nun aber $\varkappa^2$ gegeben, so entsteht die Frage, ob es zu jedem Wert von $\varkappa^2$ einen der Bedingung (3) genügenden Wert von $\omega$ gibt und ob es mehrere solche gibt.

Wir beweisen zunächst, daß durch die Differentialgleichungen (1) mit den Nebenbedingungen die Funktionen $x,y,z$ eindeutig bestimmt sind.

Es ergibt sich zunächst aus (1), daß $x^2+y^2$, $\varkappa^2x^2+z^2$ konstant sind, und aus den Nebenbedingungen folgt:
\begin{equation}
y^2=1-x^2,
\qquad
z^2=1-\varkappa^2x^2.
\tag{4}
\end{equation}
Angenommen, es existiere ein zweites System denselben Bedingungen genügender Funktionen $x_1,y_1,z_1$, so ist zunächst
\[
y_1^2=1-x_1^2,
\qquad
z_1^2=1-\varkappa^2x_1^2,
\]
und eine einfache Differentiation ergibt, daß
\[
{xy_1z_1-x_1yz\over 1-\varkappa^2x^2x_1^2}
\]
konstant ist. (Das Additionstheorem der elliptischen Integrale, aus dem man die Form dieses Ausdruckes erhält, wird unmittelbar durch Differentiation bestätigt.) Aus den Nebenbedingungen folgt aber, daß diese Konstante verschwindet, also:
\[
xy_1z_1=x_1yz,
\]
woraus man leicht schließt, indem man beiderseits quadriert:
\[
x=x_1,
\qquad y=y_1,
\qquad z=z_1.
\]

\section*{\S\,52. Die unabhängige Variable $\varkappa^2$. Lineare Differentialgleichung für $K$.}

Wir zeigen nun zunächst, daß und wie man zu einem beliebig gegebenen Wert von $\varkappa^2$ wenigstens einen der Bedingung
\begin{equation}
\varkappa^2={\tht_{10}^4\over\tht_{00}^4}
\tag{1}
\end{equation}
genügenden Wert von $\omega$ finden kann. Diese Frage wird am vollständigsten beantwortet durch Zurückführung auf eine lineare Differentialgleichung zweiter Ordnung.

Eine solche Differentialgleichung ist aber in den Formeln des \S\,23 bereits enthalten. Es folgt zunächst aus \S\,23, (14):
\[
\dd\log\varkappa^2=\pi\ii\tht_{01}^4\dd\omega,
\]
oder mit Rücksicht auf
\begin{equation}
\pi\tht_{00}^2=2K,
\qquad
\pi\tht_{10}^2=2\varkappa K,
\qquad
\pi\tht_{01}^2=2\varkappa'K
\quad[\S\,42,(3),(15)]:
\tag{2}
\end{equation}
\begin{equation}
\pi\dd(\varkappa^2)=4\ii\varkappa^2\varkappa'^2K^2\dd\omega,
\tag{3}
\end{equation}
und aus \S\,23, (19):
\[
{\dd\over\dd\omega}{1\over K^2}{\dd K\over\dd\omega}
=-{4\over\pi^2}\varkappa^2\varkappa'^2K^3.
\]
Führt man vermittelst (3) für $\omega$ die Variable $\varkappa^2$ ein, so folgt
\begin{equation}
{\dd\over\dd(\varkappa^2)}
\left(\varkappa^2\varkappa'^2{\dd K\over\dd(\varkappa^2)}\right)
={1\over4}K,
\tag{4}
\end{equation}
und dies ist die gesuchte Differentialgleichung.

Setzen wir weiter
\begin{equation}
-\ii\omega={K'\over K},
\tag{5}
\end{equation}
\begin{equation}
-\ii K^2\dd\omega=K\dd K'-K'\dd K,
\tag{6}
\end{equation}
so ergibt sich nach (3):
\begin{equation}
\varkappa^2\varkappa'^2\left(K{\dd K'\over\dd(\varkappa^2)}-K'{\dd K\over\dd(\varkappa^2)}\right)=-{\pi\over4},
\tag{7}
\end{equation}
also durch abermalige Differentiation:
\[
{1\over K}{\dd\over\dd(\varkappa^2)}
\left(\varkappa^2\varkappa'^2{\dd K\over\dd(\varkappa^2)}\right)
={1\over K'}{\dd\over\dd(\varkappa^2)}
\left(\varkappa^2\varkappa'^2{\dd K'\over\dd(\varkappa^2)}\right),
\]
woraus zu ersehen, daß $K'$ das zweite partikulare Integral der Differentialgleichung (4) ist.

Diese Differentialgleichung läßt sich durch hypergeometrische Reihen integrieren (Gauss, Disq. circa seriem infinitum Werke, Bd. III, S. 125).

Setzen wir für den Augenblick $\varkappa^2=x$, $K=y$, so lautet die Differentialgleichung (4)
\begin{equation}
x(1-x){\dd^2y\over\dd x^2}+(1-2x){\dd y\over\dd x}-{1\over4}y=0,
\tag{8}
\end{equation}
und sie geht aus der allgemeinen Gaussschen Differentialgleichung
\begin{equation}
(x-x^2){\dd^2y\over\dd x^2}+[\gamma-(\alpha+\beta+1)x]{\dd y\over\dd x}-\alpha\beta y=0
\tag{9}
\end{equation}
hervor, wenn man $\alpha=\beta={1\over2}$, $\gamma=1$ setzt. Das eine ihrer partikularen Integrale ist also
\begin{equation}
\begin{aligned}
F\left({1\over2},{1\over2},1,x\right)
&=1+\sum_{\nu=1}^{\infty}\left({1\cdot3\cdot5\cdots(2\nu-1)\over2\cdot4\cdots2\nu}\right)^2x^\nu\\
&=\sum_{\nu=0}^{\infty}{\Pi(2\nu)^2\over\Pi(\nu)^4}\left({x\over16}\right)^\nu,
\end{aligned}
\tag{10}
\end{equation}
und diese Reihe ist konvergent, so lange der absolute Wert von $x$ kleiner als $1$ ist. Als das zweite partikulare Integral kann man wählen
\begin{equation}
F\left({1\over2},{1\over2},1,1-x\right),
\tag{11}
\end{equation}
das aber einen anderen Konvergenzbereich hat. Denken wir uns $x$ als komplexe Variable in einer Ebene dargestellt, so konvergiert (10) in einem mit dem Radius 1 um den Nullpunkt beschriebenen Kreise, (11) in einem gleichen Kreise um den Punkt 1 beschrieben, so daß der gemeinsame Konvergenzbereich aus einem Zweieck besteht, das den zwischen 0 und 1 gelegenen Teil der reellen Achse enthält. Man kann aber auch, was für uns wichtig ist, das zweite partikulare Integral in einer Form aufstellen, die in demselben Kreise wie die Reihe (10) konvergiert. Dabei ist zu beachten, daß das zweite partikulare Integral für $x=0$ unendlich wird.

Dieses zweite partikulare Integral erhält man durch den folgenden Grenzübergang: Bezeichnen wir die hypergeometrische Reihe nach Gauss mit
\[
\begin{aligned}
F(\alpha,\beta,\gamma,x)
&=1+{\alpha\beta\over1\cdot\gamma}x
+{\alpha(\alpha+1)\beta(\beta+1)\over1\cdot2\cdot\gamma(\gamma+1)}x^2+\cdots\\
&={\Pi(\gamma-1)\over\Pi(\alpha-1)\Pi(\beta-1)}
\sum_{\nu=0}^{\infty}{\Pi(\alpha+\nu-1)\Pi(\beta+\nu-1)\over\Pi(\gamma+\nu-1)\Pi(\nu)}x^\nu,
\end{aligned}
\]
so erhält man im allgemeinen die beiden partikularen Integrale von (9):
\[
y_1=F_1(x)=F(\alpha,\beta,\gamma,x),
\]
\[
y_2=F_2'(x)=x^{1-\gamma}(1-x)^{\gamma-\alpha-\beta}F(1-\alpha,1-\beta,2-\gamma,x).
\]
Da diese aber für $\alpha=\beta={1\over2}$, $\gamma=1$ miteinander identisch werden, so setze man zunächst $\alpha=\beta={1\over2}$, $\gamma=1+\varepsilon$, also
\[
y_1=F\left({1\over2},{1\over2},1+\varepsilon,x\right),
\qquad
 y_2=\left({1-x\over x}\right)^{\varepsilon}F\left({1\over2},{1\over2},1-\varepsilon,x\right).
\]
Es ist aber
\[
F\left({1\over2},{1\over2},1\pm\varepsilon,x\right)
={\Pi(\pm\varepsilon)\over\Pi\left(-{1\over2}\right)^2}
\sum_{\nu=0}^{\infty}
{\Pi\left(\nu-{1\over2}\right)^2\over\Pi(\nu\pm\varepsilon)\Pi(\nu)}x^\nu,
\]
und man kann unter Hinzufügung eines konstanten Faktors als zweites partikulares Integral auch folgendes annehmen:
\[
{1\over2\varepsilon}\sum_{\nu=0}^{\infty}{\Pi\left(\nu-{1\over2}\right)^2\over\Pi(\nu)}
\left[-{1\over\Pi(\nu+\varepsilon)}+
\left({x\over1-x}\right)^{-\varepsilon}{1\over\Pi(\nu-\varepsilon)}\right]x^\nu.
\]
Entwickelt man hier nach Potenzen von $\varepsilon$ und setzt dann $\varepsilon=0$, so folgt:
\[
\sum_{\nu=0}^{\infty}{\Pi\left(\nu-{1\over2}\right)^2\over\Pi(\nu)^2}x^\nu
\left[{\Pi'(\nu)\over\Pi(\nu)}-{1\over2}\log{x\over1-x}\right],
\]
wofür man mit Benutzung der Gaussschen Formel
\[
{\Pi\left(\nu-{1\over2}\right)\Pi(\nu)\over\Pi(2\nu)}=\sqrt\pi\,2^{-2\nu}
\]
mit Abwerfung eines konstanten Faktors auch setzen kann:
\begin{equation}
\sum_{\nu=0}^{\infty}{\Pi(2\nu)^2\over\Pi(\nu)^4}\left({x\over16}\right)^\nu
\left[{\Pi'(\nu)\over\Pi(\nu)}-{1\over2}\log{x\over1-x}\right].
\tag{12}
\end{equation}
Es ist aber für ein ganzzahliges positives $\nu$:
\[
{\Pi'(\nu)\over\Pi(\nu)}=\Pi'(0)+1+{1\over2}+{1\over3}+\cdots+{1\over\nu},
\]
und wenn wir also zur Abkürzung setzen:
\begin{equation}
\begin{aligned}
F(x)&=\sum_{\nu=0}^{\infty}{\Pi(2\nu)^2\over\Pi(\nu)^4}\left({x\over16}\right)^\nu,\\
G(x)&=2\sum_{\nu=1}^{\infty}{\Pi(2\nu)^2\over\Pi(\nu)^4}\left(1+{1\over2}+{1\over3}+\cdots+{1\over\nu}\right)\left({x\over16}\right)^\nu,
\end{aligned}
\tag{13}
\end{equation}
so ergeben sich die beiden partikularen Lösungen der Differentialgleichung (8), wenn man für das zweite eine lineare Kombination von (10) und (12) nimmt:
\begin{equation}
y_1=F(x),
\qquad
y_2={1\over2}G(x)-{1\over2}\log{x\over16(1-x)}F(x).
\tag{14}
\end{equation}

Um nun aber die partikularen Integrale $K$ und $K'$ durch diese Funktionen darzustellen, müssen wir das Verhalten von $K,K'$ für $q=0$ untersuchen, wofür zugleich $\varkappa^2=0$ wird.

Es ist aber für $q=0$ nach \S\,25 und \S\,42:
\[
q^{-1}\varkappa^2=16,
\qquad
2K=\pi,
\]
also:
\[
\lim\left(\log{\varkappa^2\over16}-\ii\pi\omega\right)=0;
\]
andererseits ist $2\ii K'=\pi\tht_{00}^2\omega$ (\S\,42) oder:
\[
2K'+\log{\varkappa^2\over16}=\left(\log{\varkappa^2\over16}-\ii\pi\omega\right)+\ii\pi\omega(1-\tht_{00}^2),
\]
woraus:
\begin{equation}
\lim_{\varkappa^2=0}\left(K'+{1\over2}\log{\varkappa^2\over16}\right)=0.
\tag{15}
\end{equation}
Da nun $K$ und $K'$ partikulare Integrale von (8) sind (für $x=\varkappa^2$), so haben beide die Form:
\[
K=a_1y_1+a_2y_2,
\qquad
K'=a_1'y_1+a_2'y_2,
\]
worin $a_1,a_2,a_1',a_2'$ konstant sind. Da $K$ für $x=0$ endlich bleibt und den Wert $\frac12\pi$ erhält, so ist $a_2=0$, $a_1={1\over2}\pi$, und damit $K'+\frac12\log x$ endlich bleibe und [nach (15)] den Wert $\log4$ erhalte, muß $a_2'=1$ und $a_1'=0$ sein. Also ergibt sich, wenn man $x=\varkappa^2$, $1-x=\varkappa'^2$ setzt:
\begin{equation}
K={\pi\over2}F(\varkappa^2),
\qquad
K'={1\over2}G(\varkappa^2)-{1\over2}\log{\varkappa^2\over16\varkappa'^2}F(\varkappa^2)
\tag{16}
\end{equation}
und
\begin{equation}
q=e^{-\pi K'/K}={\varkappa^2\over16\varkappa'^2}
 e^{-{G(\varkappa^2)\over F(\varkappa^2)}},
\tag{17}
\end{equation}
worin, wie schon bemerkt, die Reihen $G(\varkappa^2)$, $F(\varkappa^2)$ konvergent sind, so lange der absolute Wert von $\varkappa^2$ ein echter Bruch ist. In der Differentialgleichung (4) liegen nun freilich die Mittel, die gefundenen Ausdrücke für $K,K',q$ über dies Konvergenzbereich hinaus fortzusetzen. Einfacher gelangt man dazu aber durch die Anwendung der Landenschen Transformation.

Wir begrenzen die Ebene der komplexen Werte $\varkappa^2$, indem wir längs der reellen Achse zwei Schnitte von 0 bis $-\infty$ und von 1 bis $\infty$ legen. In der so begrenzten Ebene sind dann alle Wurzeln aus $\varkappa^2,\varkappa'^2$ eindeutig dadurch bestimmt, daß sie für positive echt gebrochene $\varkappa^2$ reell und positiv sein sollen. Es ist nun nach den Formeln der Landenschen Transformation [\S\,44,(28)], wenn wir mit $\varkappa_1,\varkappa_1',K_1,K_1'$ die Funktionen von $\omega$ bezeichnen, die sich aus $\varkappa,\varkappa',K,K'$ ergeben, wenn $\omega$ durch $2\omega$ ersetzt wird:
\begin{equation}
\varkappa_1={1-\varkappa'\over1+\varkappa'},
\qquad
\varkappa_1'={2\sqrt{\varkappa'}\over1+\varkappa'},
\qquad
2K_1=(1+\varkappa')K,
\qquad
K_1'=(1+\varkappa')K'.
\tag{18}
\end{equation}
Wenn wir also in (17) $\omega$ durch $2\omega$ ersetzen und die Quadratwurzel ziehen, so folgt:
\begin{equation}
q={1\over4}{\varkappa_1\over\varkappa_1'}
 e^{-{1\over2}{G(\varkappa_1^2)\over F(\varkappa_1^2)}}.
\tag{19}
\end{equation}
Hierin erstreckt sich nun der Konvergenzbereich über den in der Ebene $\varkappa_1^2$ gelegenen Einheitskreis und dieser entspricht der ganzen Ebene $\varkappa^2$. Denn der reelle Teil vom $\varkappa'$ ist in der ganzen Ebene der $\varkappa^2$ positiv mit Ausnahme der beiden Ufer des von $+1$ nach $+\infty$ verlaufenden Schnittes, an denen $\varkappa'$ rein imaginär ist. Daraus ergibt sich, daß der absolute Wert von $\varkappa_1$ kleiner als 1 ist, und daß die Peripherie des Einheitskreises in der Ebene $\varkappa_1$ den beiden Ufern dieses Schnittes entspricht.

Man kann aber die Konvergenz dieser Ausdrücke noch verbessern durch eine abermalige Anwendung der Landenschen Transformation. Bezeichnen wir das Resultat einer nochmaligen Verdoppelung von $\omega$ durch den Index 2, so folgt:
\begin{equation}
\sqrt{\varkappa_2}={1-\sqrt{\varkappa'}\over1+\sqrt{\varkappa'}},
\qquad
4K_2=(1+\sqrt{\varkappa'})^2K,
\qquad
K_2'=(1+\sqrt{\varkappa'})^2K',
\qquad
q=\sqrt{{\varkappa_2\over4\varkappa_2'}e^{-{1\over4}{G(\varkappa_2^2)\over F(\varkappa_2^2)}}},
\tag{20}
\end{equation}
worin nun der Konvergenzbereich die Ebene der $\varkappa^2$ zweimal (mit einer Verzweigung im Punkte $\varkappa^2=1$) bedeckt.

Diese Operation kann man fortsetzen, indem man nach der Formel (17) aus $\varkappa_2$ durch Verdoppelung von $\omega$ eine neue Größe $\varkappa_3$, daraus ebenso $\varkappa_4$ usf. herleitet. Man bekommt dann eine Reihe von Ausdrücken für $q$, deren Konvergenz eine immer bessere wird. Für ein beliebiges $\nu$ ist:
\begin{equation}
q=\sqrt[2^{\nu-1}]{ {\varkappa_\nu\over4\varkappa_\nu'}
 e^{-{1\over2^\nu}{G(\varkappa_\nu^2)\over F(\varkappa_\nu^2)}} }.
\tag{21}
\end{equation}
Die erste Gleichung (18)
\[
\varkappa_1={\varkappa^2\over(1+\varkappa')^2}
\]
zeigt, daß, so lange der absolute Wert von $\varkappa^2$ kleiner als 1 ist, und folglich der absolute Wert von $1+\varkappa'$ größer als 1, der absolute Wert von $\varkappa_1$ kleiner ist als der von $\varkappa^2$, und folglich nähert sich $\varkappa_\nu$ mit unendlich wachsendem $\nu$ der Grenze Null. Daraus erhält man für $q$ die Darstellung:
\begin{equation}
q=\lim_{\nu=\infty}\sqrt[2^{\nu-1}]{ {\varkappa_\nu\over4\varkappa_\nu'} }
=\lim_{\nu=\infty}\sqrt[2^{\nu-1}]{ {\varkappa_\nu\over4} },
\tag{22}
\end{equation}
die bei reellen Werten von $\varkappa$ zur Berechnung geeignet ist.

Aus jeder der Formeln (21) kann man eine Reihenentwickelung von $q$ nach den steigenden Potenzen von $\varkappa_\nu$ herleiten, deren Konvergenz mit wachsendem $\nu$ zunimmt. Wendet man insbesondere die Formel (20) an, so erhält man eine Entwickelung, die von Weierstrass in den Monatsberichten der Berliner Akademie vom Jahre 1883 mitgeteilt ist, deren Konvergenzbereich die Ebene $\varkappa^2$ zweimal ausfüllt.

Die Koeffizienten dieser Entwickelung berechnet man am einfachsten aus:
\begin{equation}
\sqrt{\varkappa_2}={\tht_{10}(0,4\omega)\over\tht_{00}(0,4\omega)}
={2q+2q^9+2q^{25}+\cdots\over1+2q^4+2q^{16}+\cdots},
\tag{23}
\end{equation}
und erhält so nach der Methode der unbestimmten Koeffizienten:
\begin{equation}
q={\sqrt{\varkappa_2}\over2}
+2\left({\sqrt{\varkappa_2}\over2}\right)^5
+15\left({\sqrt{\varkappa_2}\over2}\right)^9
+150\left({\sqrt{\varkappa_2}\over2}\right)^{13}
+1707\left({\sqrt{\varkappa_2}\over2}\right)^{17}+\cdots.
\tag{24}
\end{equation}
Ebenso läßt sich nach (23) auch umgekehrt $\sqrt{\varkappa_2}$ in eine nach Potenzen von $q$ fortschreitende Reihe entwickeln:
\begin{equation}
{1\over2}\sqrt{\varkappa_2}=q-2q^5+5q^9-10q^{13}+18q^{17}\cdots,
\tag{25}
\end{equation}
die man auch durch Umkehrung der Reihe (24) erhält.

Damit ist die am Anfang dieses Paragraphen gestellte Frage vollständig beantwortet.

Bezüglich der Anwendung der hier entwickelten Formeln zu numerischer Berechnung von $q$ sei noch bemerkt, daß, wenn $\varkappa^2$ näher an 1 liegt, man ein besser konvergentes Verfahren erhält, wenn man $\varkappa^2$ mit $\varkappa'^2$ vertauscht. Die Rechnung ergibt dann zunächst nicht $q$, sondern
\[
q'=e^{-\pi\ii/\omega},
\]
woraus man aber $q$ aus der Formel erhält:
\[
\log q\log q'=\pi^2.
\]

\section*{\S\,53. Die Lösungen der Gleichung $j(\omega)=j(\omega')$.}

Die Resultate von \S\,51 genügen zunächst, um die Beziehungen der verschiedenen Werte von $\omega$ festzustellen, die zu demselben Wert von $\varkappa^2$ führen. Sei also:
\begin{equation}
{\tht_{10}^4(0,\omega')\over\tht_{00}^4(0,\omega')}
={\tht_{10}^4(0,\omega)\over\tht_{00}^4(0,\omega)}=\varkappa^2.
\tag{1}
\end{equation}
Setzen wir
\begin{equation}
v=\pi\tht_{00}^2(0,\omega)u=\pi\tht_{00}^2(0,\omega')u',
\tag{2}
\end{equation}
so genügen (nach \S\,42) sowohl die Funktionen
\begin{equation}
{\tht_{00}(0,\omega)\tht_{11}(u,\omega)\over\tht_{10}(0,\omega)\tht_{01}(u,\omega)},\quad
{\tht_{01}(0,\omega)\tht_{10}(u,\omega)\over\tht_{10}(0,\omega)\tht_{01}(u,\omega)},\quad
{\tht_{01}(0,\omega)\tht_{00}(u,\omega)\over\tht_{00}(0,\omega)\tht_{01}(u,\omega)},
\tag{3}
\end{equation}
als auch
\begin{equation}
{\tht_{00}(0,\omega')\tht_{11}(u',\omega')\over\tht_{10}(0,\omega')\tht_{01}(u',\omega')},\quad
{\tht_{01}(0,\omega')\tht_{10}(u',\omega')\over\tht_{10}(0,\omega')\tht_{01}(u',\omega')},\quad
{\tht_{01}(0,\omega')\tht_{00}(u',\omega')\over\tht_{00}(0,\omega')\tht_{01}(u',\omega')},
\tag{4}
\end{equation}
für $x,y,z$ gesetzt, den Differentialgleichungen (1) des vorigen Paragraphen, und die entsprechenden dieser Funktionen sind also identisch. Wenn nun $u'={1\over2}$ ist, so verschwindet $\tht_{10}(u',\omega')$ und es muß dann auch $\tht_{10}(u,\omega)$ verschwinden. Demnach folgt aus (2) mit Rücksicht auf \S\,21:
\begin{equation}
\tht_{00}^2(0,\omega')=\tht_{00}^2(0,\omega)(\alpha+\beta\omega),
\tag{5}
\end{equation}
worin $\alpha,\beta$ ganze Zahlen sind, die der Bedingung
\begin{equation}
\alpha\equiv1,
\qquad
\beta\equiv0\pmod 2
\tag{6}
\end{equation}
genügen. Ist $u'=\omega'/2$, so verschwinden $\tht_{01}(u',\omega')$ und $\tht_{01}(u,\omega)$, und daher ist wie oben
\begin{equation}
\tht_{00}^2(0,\omega')\omega'=\tht_{00}^2(0,\omega)(\gamma+\delta\omega),
\tag{7}
\end{equation}
worin
\begin{equation}
\gamma\equiv0,
\qquad
\delta\equiv1\pmod 2.
\tag{8}
\end{equation}
Daraus folgt:
\begin{equation}
\omega'={\gamma+\delta\omega\over\alpha+\beta\omega}.
\tag{9}
\end{equation}
Da aber in gleicher Weise geschlossen werden kann:
\[
\tht_{00}^2(0,\omega)=\tht_{00}^2(0,\omega')(\alpha'+\beta'\omega'),
\]
\[
\tht_{00}^2(0,\omega)\omega=\tht_{00}^2(0,\omega')(\gamma'+\delta'\omega'),
\]
worin $\alpha',\beta',\gamma',\delta'$ gleichfalls ganze Zahlen sind, so ergibt die Substitution (5) in diesen Gleichungen:
\[
1=(\alpha+\beta\omega)(\alpha'+\beta'\omega'),
\qquad
\omega=(\alpha+\beta\omega)(\gamma'+\delta'\omega').
\]
Drückt man in diesen Gleichungen $\omega'$ nach (9) durch $\omega$ aus, so erhält man
\[
1=\alpha'(\alpha+\beta\omega)+\beta'(\gamma+\delta\omega),
\]
\[
\omega=\gamma'(\alpha+\beta\omega)+\delta'(\gamma+\delta\omega),
\]
und da $\omega$ nicht reell ist, so zerfällt jede dieser Gleichungen in zwei andere:
\[
\alpha\alpha'+\gamma\beta'=1,
\qquad
\beta\alpha'+\delta\beta'=0,
\]
\[
\alpha\gamma'+\gamma\delta'=0,
\qquad
\beta\gamma'+\delta\delta'=1.
\]
Mithin ist
\[
(\alpha\delta-\beta\gamma)(\alpha'\delta'-\beta'\gamma')=1,
\]
und daher, da $\omega,\omega'$ beide einen positiven imaginären Teil haben müssen, also nach (9) die Determinante $\alpha\delta-\beta\gamma$ positiv sein muß:
\[
\alpha\delta-\beta\gamma=1,
\qquad
\alpha=\delta',\quad \delta=\alpha',\quad \gamma=-\gamma',\quad \beta=-\beta'.
\]

1. Es hängen also $\omega,\omega'$ durch eine lineare Transformation, und zwar [nach (6), (8)] durch eine der ersten Klasse (\S\,36) miteinander zusammen. Daß auch umgekehrt zwei solche Werte $\omega,\omega'$ denselben Wert $\varkappa^2$ ergeben, ist bereits oben nachgewiesen.

Ein ähnlicher Satz ergibt sich als unmittelbare Konsequenz hieraus, für die Invariante $j(\omega)$ (\S\,46).

2. Die Gleichung:
\begin{equation}
j(\omega')=j(\omega)
\tag{10}
\end{equation}
ist dann und nur dann befriedigt, wenn $\omega'$ mit $\omega$ durch irgend eine lineare Transformation
\[
\begin{pmatrix}\alpha&\beta\\ \gamma&\delta\end{pmatrix}
\]
zusammenhängt, wenn also
\begin{equation}
\omega'={\gamma+\delta\omega\over\alpha+\beta\omega},
\qquad
\alpha\delta-\beta\gamma=1
\tag{11}
\end{equation}
ist.

Denn bezeichnen wir die zu $\omega'$ gehörigen Werte von $\varkappa^2,\varkappa'^2$ mit $\lambda^2,\lambda'^2$, so kann die Gleichung (10) so geschrieben werden:
\[
{(1-\lambda^2\lambda'^2)^3\over\lambda^4\lambda'^4}
={ (1-\varkappa^2\varkappa'^2)^3\over\varkappa^4\varkappa'^4},
\]
und ist in bezug auf $\lambda^2$ eine Gleichung sechsten Grades, deren Wurzeln
\[
\varkappa^2,
\quad \varkappa'^2,
\quad {1\over\varkappa^2},
\quad {1\over\varkappa'^2},
\quad -{\varkappa'^2\over\varkappa^2},
\quad -{\varkappa^2\over\varkappa'^2}
\]
sind. Es findet daher (mit Rücksicht auf \S\,45) eine der folgenden Beziehungen statt:
\[
\lambda^2=\varkappa^2(\omega),
\quad \varkappa^2\left(-{1\over\omega}\right),
\quad \varkappa^2(\omega+1),
\quad \varkappa^2\left({\omega-1\over\omega}\right),
\quad \varkappa^2\left(-{1\over\omega+1}\right),
\quad \varkappa^2\left({\omega\over1-\omega}\right),
\]
wonach aus dem ersten Satze die Richtigkeit des zweiten folgt.

Die Variable $\omega$ ist hier immer auf einen positiven imaginären Teil beschränkt, und wenn wir zwei durch eine Gleichung (11) zusammenhängende Zahlen $\omega,\omega'$ äquivalent nennen, so haben alle mit $\omega$ äquivalenten Zahlen einen positiven imaginären Teil. Wir können unseren Satz 2. auch so aussprechen:

3. Die Funktion $j(\omega)$ hat für alle äquivalenten Werte $\omega$ und nur für diese einen und denselben Wert.\footnote{Dieser Satz ist zuerst von Dedekind bewiesen (Crelle, Bd. 83). Dedekind nennt danach die Funktion $\operatorname{val}(\omega)=3^{-3}2^{-6}j(\omega)$ die Valenz von $\omega$.}

Zu jedem Wert von $\omega$ gehört ein bestimmter Wert von $j(\omega)$, und zu jedem Wert von $j(\omega)$ ein Wert von $\omega$ und die Gesamtheit der äquivalenten Werte $\omega'$. Wenn also eine einwertige Funktion $\Phi(\omega)$ von $\omega$ der Bedingung genügt:
\[
\Phi\left({\gamma+\delta\omega\over\alpha+\beta\omega}\right)=\Phi(\omega),
\]
so ist $\Phi(\omega)$ eine einwertige Funktion von $j(\omega)$, und wenn wir also annehmen, daß $\Phi(\omega)$ nur für eine endliche Anzahl von Werten $j(\omega)$ unendlich und nur in endlicher Ordnung unendlich werde, so folgt, daß $\Phi(\omega)$ eine rationale Funktion von $j(\omega)$ ist.

Die Voraussetzung trifft z. B. immer dann zu, wenn $\Phi(\omega)$ mit $j(\omega)$ in einem algebraischen Zusammenhange steht.

\section*{\S\,54. Die Modulfunktionen.}

Es sei $\psi(\omega)$ eine eindeutige Funktion von $\omega$. Wendet man auf $\omega$ eine lineare Transformation
\[
S=\begin{pmatrix}\alpha&\beta\\ \gamma&\delta\end{pmatrix}
\]
an, so geht $\psi(\omega)$ in eine andere Funktion
\begin{equation}
\psi\left({\gamma+\delta\omega\over\alpha+\beta\omega}\right)
\tag{1}
\end{equation}
über, die wir mit $\psi\mid S$ bezeichnen wollen.

Die Funktion $\psi(\omega)$ kann möglicherweise bei gewissen Transformationen $S$ ungeändert bleiben (z. B. immer bei der identischen Transformation). Alle Transformationen, die eine Funktion $\psi(\omega)$ ungeändert lassen, bilden eine Gruppe $\mathfrak G$. Denn ist
\[
\psi\mid S=\psi,
\qquad
\psi\mid S_1=\psi,
\]
so ist auch
\[
\psi\mid SS_1=\psi\mid S_1=\psi.
\]
Die Gruppe $\mathfrak G$ ist ein Teiler der Gruppe $\mathfrak L$ aller linearen Transformationen.

Ist $\mathfrak G$ ein Teiler von $\mathfrak L$ von endlichem Index $(\mathfrak L,\mathfrak G)$, und $\psi(\omega)$ eine Funktion von $\omega$ von der Eigenschaft, daß auch umgekehrt
\[
\psi(\omega')=\psi(\omega)
\]
einen linearen Zusammenhang
\[
\omega'={\gamma+\delta\omega\over\alpha+\beta\omega}
\]
zur Folge hat, in dem
\[
\begin{pmatrix}\alpha&\beta\\ \gamma&\delta\end{pmatrix}
\]
eine zu $\mathfrak G$ gehörige Transformation ist, so heißt $\psi(\omega)$ eine zur Gruppe $\mathfrak G$ gehörige Modulfunktion.

Wir beschränken uns auf die Betrachtung solcher Modulfunktionen, die mit dem Modul $\varkappa^2$, also auch mit $j(\omega)$, in einem algebraischen Zusammenhange stehen.

Unter dieser Voraussetzung läßt sich der folgende allgemeine Satz aussprechen:

Wenn $\psi(\omega)$ zur Gruppe $\mathfrak G$ gehört und $\chi(\omega)$ durch die Transformationen von $\mathfrak G$ ungeändert bleibt, so ist $\chi(\omega)$ rational durch $\psi(\omega)$ ausdrückbar.

Denn zunächst ist $\chi(\omega)$ eine algebraische Funktion von $\psi(\omega)$, und $\psi(\omega)$ kann als algebraische Funktion von $j(\omega)$ jeden Wert annehmen. Zu jedem Wert von $\psi(\omega)$ gehört aber nach Voraussetzung nur ein Wert von $\chi(\omega)$, und daher ist $\chi(\omega)$ als einwertige algebraische Funktion von $\psi(\omega)$ rational.

Das Studium der in $\mathfrak L$ enthaltenen Gruppen und der zu ihnen gehörigen Modulfunktionen bildet den Hauptgegenstand des großen Werkes von Klein und Fricke: „Vorlesungen über die Theorie der elliptischen Modulfunktionen“ (2 Bde., Leipzig, Teubner, 1890, 1892). Wir betrachten hier nur einige ganz spezielle Fälle dieser Gruppen und Funktionen, auf die wir im Verlauf unserer Untersuchungen gestoßen sind.

Eine große Klasse von Teilern der Gruppe $\mathfrak L$ mit endlichem Index bilden die sogenannten Kongruenzgruppen, die, wenn $m$ irgend eine ganze Zahl ist, durch die vier Kongruenzen
\[
\begin{pmatrix}\alpha&\beta\\ \gamma&\delta\end{pmatrix}
\equiv
\begin{pmatrix}1&0\\0&1\end{pmatrix}
\pmod m
\]
charakterisiert sind. Die zu diesen Gruppen gehörigen Modulfunktionen heißen Kongruenzmoduln $m$ter Stufe. Für $m=2$ ist dies die Gruppe $\mathfrak A$ (\S\,36). Wir werden von den folgender Sätzen Gebrauch machen:

1. Der Modul $\varkappa^2$ ist nach \S\,52 eine zu dieser Gruppe gehörige Modulfunktion; da nach \S\,36 alle diese Transformationen aus den beiden
\[
\begin{pmatrix}1&0\\2&1\end{pmatrix},
\qquad
\begin{pmatrix}1&2\\0&1\end{pmatrix}
\]
zusammensetzbar sind, so können wir den Satz aussprechen:

Jede Modulfunktion, die durch die beiden Vertauschungen
\[
(\omega,\omega+2),
\qquad
\left(\omega,{\omega\over1+2\omega}\right)
\]
ungeändert bleibt, ist eine rationale Funktion von $\varkappa^2$.

2. Die Funktion $\varkappa^2\varkappa'^2$ gehört zu der aus der ersten und zweiten Klasse des \S\,36 gemeinschaftlich gebildeten Gruppe $\mathfrak A'$. Diese Gruppe läßt sich durch Wiederholung der beiden Transformationen
\[
\begin{pmatrix}1&0\\2&1\end{pmatrix},
\qquad
\begin{pmatrix}0&1\\-1&0\end{pmatrix}
\]
ableiten, und wir können daher den Satz aussprechen:

Jede Modulfunktion, die durch die beiden Vertauschungen
\[
(\omega,\omega+2),
\qquad
\left(\omega,-{1\over\omega}\right)
\]
ungeändert bleibt, ist eine rationale Funktion von $\varkappa^2\varkappa'^2$. Hieraus folgt noch:

3. Jede Modulfunktion, die durch $(\omega,\omega+2)$ ungeändert bleibt und durch $\left(\omega,-{1\over\omega}\right)$ ihr Zeichen ändert, ist das Produkt von $(\varkappa'^2-\varkappa^2)$ mit einer rationalen Funktion von $\varkappa^2\varkappa'^2$.

4. Die Invariante $j(\omega)$ gehört zur Gruppe $\mathfrak L$ (\S\,53) und daher der Satz:

Jede Modulfunktion, die durch die beiden Vertauschungen
\[
(\omega,\omega+1),
\qquad
\left(\omega,-{1\over\omega}\right)
\]
ungeändert bleibt, ist eine rationale Funktion von $j(\omega)$.

Außer diesen führen wir noch zwei andere Modulfunktionen ein. Aus der Definition der Funktionen $f(\omega),f_1(\omega),f_2(\omega)$ [\S\,34,(1),(10)] ergibt sich, wenn man die beiden Gleichungen \S\,42,(3) miteinander multipliziert:
\begin{equation}
f(\omega)={\sqrt[16]{2}\over\sqrt[12]{\varkappa\varkappa'}},
\qquad
f_1(\omega)=\sqrt[16]{2}\sqrt[12]{\varkappa'^2\over\varkappa},
\qquad
f_2(\omega)=\sqrt[16]{2}\sqrt[12]{\varkappa^2\over\varkappa'},
\tag{2}
\end{equation}
\begin{equation}
{f_1(\omega)\over f(\omega)}=\sqrt[4]{\varkappa'},
\qquad
{f_2(\omega)\over f(\omega)}=\sqrt[4]{\varkappa}.
\tag{3}
\end{equation}
Auf Grund von \S\,46,(15),(16) definieren wir zwei Funktionen $\gamma_2(\omega),\gamma_3(\omega)$:
\begin{equation}
\begin{aligned}
\gamma_2(\omega)&=\sqrt[3]{j(\omega)}
=\sqrt[3]{2^8}\,{1-\varkappa^2\varkappa'^2\over\sqrt[3]{\varkappa^4\varkappa'^4}},\\
\gamma_3(\omega)&=\sqrt{j(\omega)-27\cdot64}
={8(2+\varkappa^2\varkappa'^2)(\varkappa'^2-\varkappa^2)\over\varkappa^2\varkappa'^2},
\end{aligned}
\tag{4}
\end{equation}
die nach (2) und (3) als eindeutige Funktionen von $\omega$ folgendermaßen darstellbar sind:
\begin{equation}
\gamma_2(\omega)={f(\omega)^{24}-16\over f(\omega)^8},
\qquad
\gamma_3(\omega)={\bigl[f(\omega)^{24}+8\bigr]\bigl[f_1(\omega)^8-f_2(\omega)^8\bigr]\over f(\omega)^8},
\tag{5}
\end{equation}
wofür man nach den Grundformeln für die $f$-Funktionen [\S\,34,(11),(12)] auch setzen kann:
\begin{equation}
\gamma_2(\omega)=f(\omega)^8f_1(\omega)^8+f(\omega)^8f_2(\omega)^8-f_1(\omega)^8f_2(\omega)^8,
\tag{6}
\end{equation}
\begin{equation}
\gamma_3(\omega)={1\over2}\bigl[f(\omega)^8-f_1(\omega)^8\bigr]
\bigl[f(\omega)^8+f_2(\omega)^8\bigr]
\bigl[f_1(\omega)^8-f_2(\omega)^8\bigr].
\tag{7}
\end{equation}
Es sind also $x=f^8,-f_1^8,-f_2^8$ die Wurzeln der kubischen Gleichung
\begin{equation}
x^3-\gamma_2x-16=0,
\tag{8}
\end{equation}
und $4\gamma_3^2$ ist die Diskriminante dieser Gleichung.

Bemerkenswert sind auch die Differentialgleichungen:
\begin{equation}
\dd\varkappa^2=-\dd\varkappa'^2=\pi\ii\tht_{00}^4\varkappa^2\varkappa'^2\dd\omega
\quad[\S\,52,(3)],
\tag{9}
\end{equation}
\begin{equation}
\dd\varkappa^2\varkappa'^2=-\pi\ii\tht_{00}^4(\varkappa^2-\varkappa'^2)\varkappa^2\varkappa'^2\dd\omega,
\tag{10}
\end{equation}
\begin{equation}
\dd f(\omega)^8={\pi\ii\over3}\tht_{00}^4\bigl[f_2(\omega)^8-f_1(\omega)^8\bigr]\dd\omega,
\tag{11}
\end{equation}
\begin{equation}
\dd\gamma_2(\omega)=-{2\pi\ii\over3}\gamma_3(\omega)\eta(\omega)^4\dd\omega
\quad\text{[nach (5) und \S\,34,(10)]},
\tag{12}
\end{equation}
\begin{equation}
\dd j(\omega)=-2\pi\ii\gamma_2(\omega)^2\gamma_3(\omega)\eta(\omega)^4\dd\omega.
\tag{13}
\end{equation}
Man erhält sodann für die linearen Fundamentaltransformationen aus \S\,34,(13),(14),(15):
\begin{equation}
\gamma_2(\omega+1)=e^{-{2\pi\ii\over3}}\gamma_2(\omega),
\qquad
\gamma_2\left(-{1\over\omega}\right)=\gamma_2(\omega),
\tag{14a}
\end{equation}
\begin{equation}
\gamma_3(\omega+1)=-\gamma_3(\omega),
\qquad
\gamma_3\left(-{1\over\omega}\right)=-\gamma_3(\omega),
\tag{14b}
\end{equation}
und aus \S\,40 findet man allgemein:
\begin{equation}
\gamma_2\left({\gamma+\delta\omega\over\alpha+\beta\omega}\right)=\varrho\gamma_2(\omega),
\qquad
\gamma_3\left({\gamma+\delta\omega\over\alpha+\beta\omega}\right)=(-1)^{\alpha\gamma+\beta\gamma+\beta\delta}\gamma_3(\omega),
\tag{15}
\end{equation}
worin $\varrho$ wie in \S\,40,(3) die Bedeutung hat:
\[
\varrho=e^{-{2\pi\ii\over3}(\alpha\gamma+\beta\delta-\alpha\beta-\alpha^2\beta\delta)}.
\]
Es sind dies also Modulfunktionen, und die Gruppen, zu denen sie gehören, sind durch die beiden Kongruenzen charakterisiert:
\[
-\alpha\beta+\alpha\gamma+\beta\delta-\alpha^2\beta\gamma\equiv0\pmod3,
\]
\[
\alpha\gamma+\beta\gamma+\beta\delta\equiv0\pmod2.
\]
Wir geben hiernach unserem Satz 4. die folgende Ergänzung, die sich aus den Formeln (14) ergibt.

5. Eine Modulfunktion, die durch die beiden Substitutionen
\[
(\omega,\omega+1),
\qquad
\left(\omega,-{1\over\omega}\right),
\]
das Zeichen ändert, ist das Produkt von $\gamma_3(\omega)$ mit einer rationalen Funktion von $j(\omega)$.

6. Eine Modulfunktion, die durch die Substitution $\left(\omega,-{1\over\omega}\right)$ ungeändert bleibt und durch $(\omega,\omega+1)$ den Faktor $e^{\pm2\pi\ii/3}$ annimmt, ist das Produkt von $\gamma_2(\omega)^{\pm1}$ mit einer rationalen Funktion $j(\omega)$.

7. Eine Modulfunktion, die durch die Substitutionen $\left(\omega,-{1\over\omega}\right)$ das Zeichen ändert und durch $(\omega,\omega+1)$ den Faktor $-e^{\pm2\pi\ii/3}$ annimmt, ist das Produkt von $\gamma_2(\omega)\gamma_3(\omega)^{\pm1}$ mit einer rationalen Funktion von $j(\omega)$.

Eingehender werden wir uns mit den Modulfunktionen im zweiten Teil beschäftigen.

\section*{\S\,55. Darstellung der elliptischen Funktionen durch $v$ und $\varkappa^2$.}

Wenn wir das Umkehrproblem der elliptischen Integrale so wie in \S\,51 als die Aufgabe der Integration der elliptischen Differentialgleichungen fassen, so verlangt es die Darstellung der drei Funktionen $\sn v,\cn v,\dn v$ durch die beiden unabhängigen Variablen $v,\varkappa^2$. Diese Aufgabe ist durch \S\,42 gelöst, aber bezüglich der Variablen $\varkappa^2$ nur implizite. Man kann aber auch Darstellungen finden, in denen $\varkappa^2$ explizite vorkommt, und zwar durch Reihen, die nach Potenzen von $v$ fortschreiten, deren Koeffizienten rationale Funktionen von $\varkappa^2$ sind. Die Koeffizienten dieser Reihen können freilich nicht durch ein allgemeines Gesetz dargestellt, sondern nur sukzessive berechnet werden. Diese Darstellungen verdankt man hauptsächlich Weierstrass.\footnote{Über die Weierstrasssche Theorie der elliptischen Funktionen vgl. man besonders die von H. A. Schwarz herausgegebenen „Formeln und Lehrsätze zum Gebrauch der elliptischen Funktionen“. Über die Reihenentwickelungen vgl. man auch Königsberger, „Vorlesungen über die Theorie der elliptischen Funktionen“, 25. Vorlesung.}

Die $\sigma$-Funktionen können, da sie durchaus endliche und stetige Funktionen von $u$ sind, in unbedingt konvergente Potenzreihen nach $u$ entwickelt werden, und wenn wir daher die in \S\,45,(2) vorkommenden $\sigma$-Funktionen in dieser Weise entwickeln, so bekommen wir eine Lösung unserer Aufgabe. Wir setzen daher
\begin{equation}
\begin{aligned}
\sigma(v,2K,2\ii K')&=\sum_{\nu=0}^{\infty}A_\nu{v^{2\nu+1}\over(2\nu+1)!},\\
\sigma_{00}(v,2K,2\ii K')&=\sum_{\nu=0}^{\infty}B_\nu{v^{2\nu}\over(2\nu)!},\\
\sigma_{01}(v,2K,2\ii K')&=\sum_{\nu=0}^{\infty}C_\nu{v^{2\nu}\over(2\nu)!},\\
\sigma_{10}(v,2K,2\ii K')&=\sum_{\nu=0}^{\infty}D_\nu{v^{2\nu}\over(2\nu)!},
\end{aligned}
\tag{1}
\end{equation}
und die $A_\nu,B_\nu,C_\nu,D_\nu$ sind die zu bestimmenden Funktionen von $\varkappa^2$ oder von $\omega$. Fassen wir sie zunächst als Funktionen von $\omega$ auf, so ergeben die linearen Transformationen [\S\,45,(3)]:
\begin{equation}
\begin{aligned}
A_\nu\left(-{1\over\omega}\right)&=(-1)^\nu A_\nu(\omega),\\
B_\nu\left(-{1\over\omega}\right)&=(-1)^\nu B_\nu(\omega),\\
C_\nu\left(-{1\over\omega}\right)&=(-1)^\nu D_\nu(\omega),\\
D_\nu\left(-{1\over\omega}\right)&=(-1)^\nu C_\nu(\omega),
\end{aligned}
\tag{2}
\end{equation}
\begin{equation}
\begin{aligned}
\varkappa'^{2\nu}A_\nu(\omega+1)&=A_\nu(\omega),\\
\varkappa'^{2\nu}B_\nu(\omega+1)&=C_\nu(\omega),\\
\varkappa'^{2\nu}C_\nu(\omega+1)&=B_\nu(\omega),\\
\varkappa'^{2\nu}D_\nu(\omega+1)&=D_\nu(\omega),
\end{aligned}
\tag{3}
\end{equation}
\begin{equation}
A_\nu(\omega+2)=A_\nu(\omega),
\quad
B_\nu(\omega+2)=B_\nu(\omega),
\quad
C_\nu(\omega+2)=C_\nu(\omega),
\quad
D_\nu(\omega+2)=D_\nu(\omega).
\tag{4}
\end{equation}
Hieraus schließen wir auf die charakteristischen Eigenschaften dieser Koeffizienten als Funktionen von $\varkappa^2$. Zunächst erhellt aus der Bedeutung der Koeffizienten, daß für jeden Wert von $\varkappa^2$ mit etwaiger Ausnahme von $\varkappa^2=0,1,\infty$ die Koeffizienten $A_\nu(\varkappa^2),B_\nu(\varkappa^2),C_\nu(\varkappa^2),D_\nu(\varkappa^2)$ endliche und stetige Funktionen von $\varkappa^2$ sind.

Aber auch für $\varkappa^2=0$ bleiben diese Funktionen endlich und man kann ihre Werte leicht finden, wenn man in den Darstellungen des \S\,37 $q$ und mithin $\varkappa^2$ in Null übergehen läßt. Man erhält dann aus den Entwickelungen in \S\,25 mit Rücksicht darauf, daß nach \S\,34 und \S\,42 für ein verschwindendes $q$
\[
\eta(\omega)=q^{1/12},
\qquad
2K=\pi
\]
wird:
\begin{equation}
\begin{aligned}
e^{v^2/6}\sin v&=\sum_{\nu=0}^{\infty}A_\nu(0){v^{2\nu+1}\over(2\nu+1)!},\\
e^{v^2/6}&=\sum_{\nu=0}^{\infty}B_\nu(0){v^{2\nu}\over(2\nu)!},\\
e^{v^2/6}&=\sum_{\nu=0}^{\infty}C_\nu(0){v^{2\nu}\over(2\nu)!},\\
e^{v^2/6}\cos v&=\sum_{\nu=0}^{\infty}D_\nu(0){v^{2\nu}\over(2\nu)!}.
\end{aligned}
\tag{5}
\end{equation}
Die Formeln (2) ergeben nun:
\begin{equation}
\begin{aligned}
A_\nu(\varkappa'^2)&=(-1)^\nu A_\nu(\varkappa^2),\\
B_\nu(\varkappa'^2)&=(-1)^\nu B_\nu(\varkappa^2),\\
C_\nu(\varkappa'^2)&=(-1)^\nu D_\nu(\varkappa^2),\\
D_\nu(\varkappa'^2)&=(-1)^\nu C_\nu(\varkappa^2),
\end{aligned}
\tag{6}
\end{equation}
woraus folgt, daß auch für $\varkappa^2=1$ diese Funktionen endlich bleiben, nämlich:
\begin{equation}
\begin{aligned}
A_\nu(1)&=(-1)^\nu A_\nu(0),\\
B_\nu(1)&=(-1)^\nu B_\nu(0),\\
C_\nu(1)&=(-1)^\nu D_\nu(0),\\
D_\nu(1)&=(-1)^\nu C_\nu(0).
\end{aligned}
\tag{7}
\end{equation}
Das System (3) läßt sich ferner so schreiben:
\begin{equation}
\begin{aligned}
\varkappa'^{2\nu}A_\nu\left(-{\varkappa^2\over\varkappa'^2}\right)&=A_\nu(\varkappa^2),\\
\varkappa'^{2\nu}B_\nu\left(-{\varkappa^2\over\varkappa'^2}\right)&=C_\nu(\varkappa^2),\\
\varkappa'^{2\nu}C_\nu\left(-{\varkappa^2\over\varkappa'^2}\right)&=B_\nu(\varkappa^2),\\
\varkappa'^{2\nu}D_\nu\left(-{\varkappa^2\over\varkappa'^2}\right)&=D_\nu(\varkappa^2),
\end{aligned}
\tag{8}
\end{equation}
woraus für ein unendliches $\varkappa^2$:
\begin{equation}
\begin{aligned}
A_\nu(\varkappa^2)&=\varkappa^{2\nu}A_\nu(0),\\
B_\nu(\varkappa^2)&=\varkappa^{2\nu}D_\nu(0),\\
C_\nu(\varkappa^2)&=\varkappa^{2\nu}B_\nu(0),\\
D_\nu(\varkappa^2)&=\varkappa^{2\nu}C_\nu(0),
\end{aligned}
\tag{9}
\end{equation}
und hieraus wird geschlossen, daß die Koeffizienten $A_\nu,B_\nu,C_\nu,D_\nu$ ganze rationale Funktionen von $\varkappa^2$ vom Grade $\nu$ sind, und aus \S\,54, 2, 3, folgt überdies noch, daß $A_\nu,B_\nu$ bei geradem $\nu$ rational durch $\varkappa^2\varkappa'^2$ ausgedrückt sind, während sie bei ungeradem $\nu$ gleich dem Produkt von $(\varkappa'^2-\varkappa^2)$ mit einer rationalen Funktion von $\varkappa^2\varkappa'^2$ sind.

Aus $B_\nu$ findet man $C_\nu$ mittels der Formel:
\begin{equation}
C_\nu(\varkappa^2)=\varkappa'^{2\nu}B_\nu\left(-{\varkappa^2\over\varkappa'^2}\right),
\tag{10}
\end{equation}
und $D_\nu$ durch Anwendung von (6) und (8):
\begin{equation}
D_\nu(\varkappa^2)=(-1)^\nu C_\nu(\varkappa'^2)=(-1)^\nu\varkappa^{2\nu}B_\nu\left(-{\varkappa'^2\over\varkappa^2}\right).
\tag{11}
\end{equation}
Da man aus (5) die Koeffizienten $A_\nu(0),B_\nu(0),C_\nu(0),D_\nu(0)$ leicht bestimmen kann, so ergeben sich aus dem hier entwickelten Formelsystem die Koeffizienten $A_\nu,B_\nu,C_\nu,D_\nu$ ohne weitere Rechnung bis $\nu=3$ einschließlich.

Man findet:
\[
A_0=1,
\quad A_1=0,
\quad A_2=-{2\over3}(1-\varkappa^2\varkappa'^2),
\]
\[
B_0=1,
\quad B_1={1\over3}(\varkappa'^2-\varkappa^2),
\quad B_2={1\over3}(1+2\varkappa^2\varkappa'^2),
\]
\[
C_0=1,
\quad C_1={1\over3}(1+\varkappa^2),
\quad C_2={1\over3}(\varkappa'^4-2\varkappa^2),
\]
\[
D_0=1,
\quad D_1=-{1\over3}(1+\varkappa'^2),
\quad D_2={1\over3}(\varkappa^4-2\varkappa'^2),
\]
\[
A_3=-{8\over9}(\varkappa'^2-\varkappa^2)(2+\varkappa^2\varkappa'^2),
\]
\[
B_3={1\over9}(\varkappa'^2-\varkappa^2)(5-2\varkappa^2\varkappa'^2),
\]
\[
C_3={1\over9}(1+\varkappa^2)(5\varkappa'^4+2\varkappa^2),
\]
\[
D_3=-{1\over9}(1+\varkappa'^2)(5\varkappa^4+2\varkappa'^2).
\]
Weitere Koeffizienten lassen sich auf demselben Wege, wenn auch weitläufiger berechnen, indem man die Ausdrücke für die $\sigma$-Funktionen in \S\,32 nach Potenzen von $q$ entwickelt und, wie hier die niedrigste Potenz von $q$, so die höheren benutzt. Weierstrass bedient sich zur rekurrenten Berechnung der Koeffizienten gewisser partieller Differentialgleichungen, ähnlich der, die wir im nächsten Paragraphen für die Funktion $\sigma$ ableiten werden.

\section*{\S\,56. Potenzreihen für die Weierstrassschen Funktionen $\wp(u),\sigma(u)$.}

Die Funktion $\sigma(u)$, als Funktion von $u$ betrachtet, läßt sich, wie schon bemerkt, in eine in der ganzen $u$-Ebene konvergierende Reihe nach Potenzen von $u$ entwickeln, und nach \S\,35,(8),(12) hat diese Entwickelung die folgende Form:
\begin{equation}
\sigma(u)=u+\alpha u^3+\alpha_1u^5+\alpha_2u^7+\alpha_3u^9+\cdots.
\tag{1}
\end{equation}
Die Funktion
\begin{equation}
\wp(u)=-{\dd^2\log\sigma(u)\over\dd u^2}
={\sigma'(u)^2-\sigma(u)\sigma''(u)\over\sigma(u)^2}
\tag{2}
\end{equation}
läßt sich ebenfalls nach steigenden Potenzen von $u$ entwickeln, aber nicht mehr in eine immer konvergente Reihe, sondern diese Reihe hat einen Konvergenzkreis. Sie hat die Form:
\begin{equation}
\wp(u)={1\over u^2}+a+a_1u^2+a_2u^4+a_3u^6+\cdots,
\tag{3}
\end{equation}
und man kann die Koeffizienten $a,a_1,a_2,a_3,\ldots$ der Reihe nach aus der Differentialgleichung [\S\,46,(18)]
\begin{equation}
\wp'(u)^2=4\wp(u)^3-g_2\wp(u)-g_3
\tag{4}
\end{equation}
als Funktionen von $g_2,g_3$ bestimmen.

Es ist zunächst
\[
\wp'(u)=-{2\over u^3}+2a_1u+4a_2u^3+6a_3u^5+\cdots,
\]
\[
\wp'(u)^2={4\over u^6}-{8a_1\over u^2}-16a_2+(4a_1^2-24a_3)u^2+\cdots,
\]
und da $\wp'(u)^2$ kein Glied mit $u^{-4}$ enthält, während in $\wp(u)^3$ das Glied $3au^{-4}$ vorkommt, so muß $a=0$ sein. Danach ergibt sich:
\[
\wp(u)^3={1\over u^6}+{3a_1\over u^2}+3a_2+3(a_3+a_1^2)u^2+\cdots,
\]
\[
\wp(u)={1\over u^2}+a_1u^2+\cdots.
\]
Und wenn man dies in die Gleichung (4) einsetzt, so folgt
\begin{equation}
a_1={g_2\over20},
\qquad
a_2={g_3\over28},
\qquad
a_3={g_2^2\over1200}.
\tag{5}
\end{equation}
Durch Integration von (2) findet man:
\[
\sigma'(u)=\left({1\over u}-{a_1u^3\over3}-{a_2u^5\over5}-{a_3u^7\over7}\cdots\right)\sigma(u),
\]
und aus dieser Gleichung folgt:
\begin{equation}
\alpha=0,
\qquad
\alpha_1=-{g_2\over240},
\qquad
\alpha_2=-{g_3\over840},
\qquad
\alpha_3=-{g_2^2\over161280}.
\tag{6}
\end{equation}
Daß $\alpha=0$ ist, folgt auch aus dem in \S\,35 bewiesenen $\sigma'''(u)=0$; indessen war dies dort auf ziemlich umständlichem Wege bewiesen. Wir haben sonach folgenden Anfang der Entwickelung von $\sigma(u)$, wobei der Übersicht halber die numerischen Nenner in ihre Primfaktoren zerlegt sind:
\begin{equation}
\sigma(u)=u-{g_2\over2^4\cdot3\cdot5}u^5
-{g_3\over2^3\cdot3\cdot5\cdot7}u^7
-{g_2^2\over2^9\cdot3^2\cdot5\cdot7}u^9-\cdots.
\tag{7}
\end{equation}

Zur rekurrenten Berechnung der Koeffizienten in der Potenzreihe für $\sigma(u)$ dient eine partielle Differentialgleichung, die als eine Umformung der partiellen Differentialgleichung für die $\vartheta$-Funktionen betrachtet werden kann. Wir leiten sie auf folgendem Wege her.

Die Funktion $\sigma(u)$ war durch die beiden Periodengleichungen \S\,35,(9)
\begin{equation}
\sigma(u+\omega_1)=-e^{\eta_1(2u+\omega_1)}\sigma(u),
\qquad
\sigma(u+\omega_2)=-e^{\eta_2(2u+\omega_2)}\sigma(u),
\qquad
\eta_1\omega_2-\eta_2\omega_1=\pi\ii
\tag{8}
\end{equation}
bis auf einen von $u$ unabhängigen Faktor bestimmt. Es ergibt sich aber durch Differentiation der letzten Gleichung (8):
\[
{\partial\eta_1\over\partial\omega_1}\omega_2
-{\partial\eta_2\over\partial\omega_1}\omega_1-
\eta_2=0,
\]
\[
{\partial\eta_1\over\partial\omega_2}\omega_2
-{\partial\eta_2\over\partial\omega_2}\omega_1+
\eta_1=0,
\]
also:
\[
\omega_2\left(\eta_1{\partial\eta_1\over\partial\omega_1}+\eta_2{\partial\eta_1\over\partial\omega_2}\right)
=\omega_1\left(\eta_1{\partial\eta_2\over\partial\omega_1}+\eta_2{\partial\eta_2\over\partial\omega_2}\right),
\]
und wir führen also eine Größe $r$ ein, die wir so definieren:
\begin{equation}
\eta_1{\partial\eta_1\over\partial\omega_1}+\eta_2{\partial\eta_1\over\partial\omega_2}=-r\omega_1,
\qquad
\eta_1{\partial\eta_2\over\partial\omega_1}+\eta_2{\partial\eta_2\over\partial\omega_2}=-r\omega_2.
\tag{9}
\end{equation}
Nun weist man durch Differentiation der Periodengleichungen (8) nach, daß die Funktion
\[
S(u)={\partial^2\sigma\over\partial u^2}
+4\eta_1{\partial\sigma\over\partial\omega_1}
+4\eta_2{\partial\sigma\over\partial\omega_2}
+4ru^2\sigma
\]
den Periodengleichungen (8) genügt, und also gleich $C\sigma(u)$ ist, worin $C$ von $u$ unabhängig ist. Die Entwickelung von $S(u)$ nach Potenzen von $u$ fängt aber mit $u^3$ an, während $\sigma(u)$ mit $u$ anfängt, und folglich muß $C=0$ sein.

Nach (7) beginnt die Entwickelung von $\partial\sigma/\partial\omega_1$ und $\partial\sigma/\partial\omega_2$ erst mit $u^5$, und die ersten Glieder von $\partial^2\sigma/\partial u^2$ und $4ru^2\sigma$ sind $-g_2u^3/12$, $4ru^3$, woraus $4r=g_2/12$ folgt, und demnach erhält die Differentialgleichung für $\sigma$ die Form:
\begin{equation}
{\partial^2\sigma\over\partial u^2}
+4\eta_1{\partial\sigma\over\partial\omega_1}
+4\eta_2{\partial\sigma\over\partial\omega_2}
+{g_2\over12}u^2\sigma=0,
\tag{10}
\end{equation}
und nebenbei ergeben sich aus (9) die Relationen:
\begin{equation}
\eta_1{\partial\eta_1\over\partial\omega_1}+\eta_2{\partial\eta_1\over\partial\omega_2}=-{g_2\over48}\omega_1,
\qquad
\eta_1{\partial\eta_2\over\partial\omega_1}+\eta_2{\partial\eta_2\over\partial\omega_2}=-{g_2\over48}\omega_2.
\tag{11}
\end{equation}

Nun sollen in der Differentialgleichung (10) statt der Variablen $\omega_1,\omega_2$ die Variablen $g_2,g_3$ eingeführt werden, und wenn wir also
\[
{\partial\sigma\over\partial\omega_1}={\partial\sigma\over\partial g_2}{\partial g_2\over\partial\omega_1}+{\partial\sigma\over\partial g_3}{\partial g_3\over\partial\omega_1},
\qquad
{\partial\sigma\over\partial\omega_2}={\partial\sigma\over\partial g_2}{\partial g_2\over\partial\omega_2}+{\partial\sigma\over\partial g_3}{\partial g_3\over\partial\omega_2},
\]
setzen, so erhält (10) die Form:
\begin{equation}
{\partial^2\sigma\over\partial u^2}
+h_2{\partial\sigma\over\partial g_2}
+h_3{\partial\sigma\over\partial g_3}
+{g_2\over12}u^2\sigma=0,
\tag{12}
\end{equation}
worin
\[
h_2=4\left(\eta_1{\partial g_2\over\partial\omega_1}+\eta_2{\partial g_2\over\partial\omega_2}\right),
\qquad
h_3=4\left(\eta_1{\partial g_3\over\partial\omega_1}+\eta_2{\partial g_3\over\partial\omega_2}\right).
\]
Die Größen $h_2,h_3$, durch $g_2,g_3$ ausgedrückt, ergeben sich nun aus der Differentialgleichung selbst, wenn man für $\sigma$ die Entwickelung (7) einsetzt. Danach ist bis zur 7ten Potenz von $u$:
\[
{\partial^2\sigma\over\partial u^2}=-{g_2\over12}u^3
-12g_3{u^5\over2^4\cdot3\cdot5}
-{3\over8}g_2^2{u^7\over2^3\cdot3\cdot5\cdot7},
\]
\[
{\partial\sigma\over\partial g_2}=-{u^5\over2^4\cdot3\cdot5},
\qquad
{\partial\sigma\over\partial g_3}=-{u^7\over2^3\cdot3\cdot5\cdot7},
\]
\[
{g_2\over12}u^2\sigma={g_2\over12}u^3
-{7\over24}g_2^2{u^7\over2^3\cdot3\cdot5\cdot7},
\]
und wenn man dies in (12) einsetzt und die Koeffizienten von $u^5$ und $u^7$ gleich Null setzt, so ergibt sich:
\[
h_2=-12g_3,
\qquad
h_3=-{2\over3}g_2^2,
\]
und wir erhalten die Differentialgleichung:
\begin{equation}
{\partial^2\sigma\over\partial u^2}
-12g_3{\partial\sigma\over\partial g_2}
-{2\over3}g_2^2{\partial\sigma\over\partial g_3}
+{g_2\over12}u^2\sigma=0.
\tag{13}
\end{equation}
Setzt man die Reihenentwickelung für $\sigma$ mit unbestimmten Koeffizienten an:
\begin{equation}
\sigma=\sum A_\nu{u^{2\nu+1}\over(2\nu+1)!},
\tag{14}
\end{equation}
so erhält man aus (13) die Rekursionsformel:
\begin{equation}
A_\nu=12g_3{\partial A_{\nu-1}\over\partial g_2}
+{2\over3}g_2^2{\partial A_{\nu-1}\over\partial g_3}
-{(2\nu-1)(2\nu-2)\over12}g_2A_{\nu-2},
\tag{15}
\end{equation}
woraus zu schließen ist, daß die $A_\nu$ ganze rationale Funktionen von $g_2$ und $g_3$ sind, deren Koeffizienten rationale Zahlen sind, die nur Potenzen von 2 und von 3 im Nenner haben können. Die Reihe (14) besteht daher aus Gliedern von der Form
\[
a_{m,n}\left({1\over2}g_2\right)^m(2g_3)^n{u^{2\nu+1}\over(2\nu+1)!},
\]
und aus der Homogenität der Funktion $\sigma$ [\S\,37,(8), \S\,46,(13)] folgt:
\begin{equation}
2m+3n=\nu,
\tag{16}
\end{equation}
für die Koeffizienten $a_{m,n}$ ergibt sich dann aus (15):
\begin{equation}
\begin{aligned}
a_{m,n}={}&3(m+1)a_{m+1,n-1}+{16\over3}(n+1)a_{m-2,n+1}\\
&-{1\over3}(2m+3n-1)(4m+6n-1)a_{m-1,n},
\end{aligned}
\tag{17}
\end{equation}
worin $a_{0,0}=1$ und die $a$, bei denen einer der beiden Indizes negativ ist, gleich Null zu setzen sind.

Daraus läßt sich $a_{m,n}$ finden, wenn die früheren $a_{m,n}$, d. h. die, in denen $2m+3n$ kleiner als $\nu$ ist, schon berechnet sind. Man sieht daraus, daß die $a_{m,n}$ nur Potenzen von 3 im Nenner enthalten. In den von H. A. Schwarz herausgegebenen „Formeln und Lehrsätzen zum Gebrauch der elliptischen Funktionen“, nach Vorlesungen von Weierstrass (2. Ausgabe 1893), sind die $a_{m,n}$ bis zu $\nu=17$ berechnet. Es zeigt sich dabei, daß diese Koeffizienten nicht nur ganze Zahlen sind, sondern daß sie auch mit $\nu$ wachsende Potenzen von 3 als Faktoren enthalten. Einen Beweis hierfür vermag ich nicht zu geben.

\end{document}
